\chapter{Obstructive Sleep Apnea}
\index{Obstructive Sleep Apnea}
\label{sec:Obstructive Sleep Apnea}

\section{Overview}
Obstructive sleep apnea (OSA) is a sleep-related breathing disorder characterized by recurrent episodes of partial or complete pharyngeal collapse during sleep, resulting in apnea (breathing cessation $\ge$10 seconds) or hypopnea (reduced airflow with desaturation or arousal), causing intermittent hypoxia, hypercapnia, intrathoracic pressure swings, and sleep fragmentation. Prevalence is 15--30\% of men and 10--20\% of women in the general population, with marked increases in obesity. This condition is among the most common mental health disorders worldwide. It affects people across all age groups, socioeconomic backgrounds, and geographic regions. This condition affects a significant number of people worldwide and can range from mild to severe in presentation. Early recognition and appropriate management are essential for optimal outcomes. A thorough understanding of the underlying mechanisms guides treatment decisions. Patient education and self-management play important roles in long-term care.
\section{Causes}
This condition happens because of anatomic narrowing of the pharyngeal airway (obesity, craniofacial abnormalities, macroglossia, tonsillar enlargement) combined with reduced dilator muscle activity of the genioglossus during sleep and a low arousal threshold. The underlying mechanisms involve complex interactions between genetic predisposition and environmental factors that disrupt normal physiological processes. The underlying mechanisms involve complex interactions between genetic predisposition and environmental factors. Disruption of normal physiological processes leads to the characteristic features of the condition. Multiple contributing factors may act synergistically to trigger or worsen the condition. Understanding the specific cause helps guide targeted treatment approaches.
\section{Risk Factors}
Genetic predisposition and family history Advancing age Lifestyle factors (diet, exercise, smoking, alcohol) Environmental exposures Underlying medical conditions Medication use Stress and psychological factors Socioeconomic factors

\begin{itemize}[noitemsep]
  \item Genetic predisposition and family history
  \item Advancing age
  \item Lifestyle factors (diet, exercise, smoking, alcohol)
  \item Environmental exposures
  \item Underlying medical conditions
  \item Medication use
\end{itemize}
\section{Signs and Symptoms}
You may experience loud snoring, witnessed apneas, gasping or choking awakenings, nocturia, morning headache, excessive daytime sleepiness, fatigue, thinking and memory impairment, and irritability. Symptoms vary depending on the severity and extent of the condition Pain or discomfort in the affected area Functional impairment affecting daily activities Changes in appearance or function of affected tissues Systemic symptoms may include fatigue, fever, or weight changes Symptoms may develop gradually or appear suddenly

\begin{itemize}[noitemsep]
  \item Pain or discomfort in the affected area
  \item Functional impairment affecting daily activities
  \item Changes in appearance or function of affected tissues
  \item Systemic symptoms may include fatigue, fever, or weight changes
  \item Symptoms may develop gradually or appear suddenly
\end{itemize}
\section{Progression and Stages}
Doctors diagnose this using attended in-laboratory polysomnography (gold standard) or home sleep apnea testing (HSAT) in uncomplicated adults, with severity classified by the apnea-hypopnea index (AHI): mild (AHI 5--15/hr), moderate (AHI 15--30/hr), and severe (AHI >30/hr). The condition may progress through different stages of severity over time. Early stages often involve mild or subtle symptoms that may be overlooked. Moderate stages involve more noticeable symptoms affecting daily function. Advanced stages may involve significant impairment and complications.
\section{Complications}
\begin{itemize}[noitemsep]
  \item Disease progression leading to worsening symptoms
  \item Secondary infections or injuries
  \item Side effects of treatment
  \item Reduced quality of life
  \item Emotional and psychological impact
\end{itemize}
\section{Diagnosis}
Comprehensive medical history and physical examination Laboratory tests to assess organ function and rule out other conditions Imaging studies to visualize affected structures Specialized testing depending on the affected system Consultation with relevant specialists Monitoring of disease progression over time

\begin{itemize}[noitemsep]
  \item Comprehensive medical history and physical examination
  \item Laboratory tests to assess organ function
  \item Imaging studies to visualize affected structures
  \item Specialized testing depending on the affected system
  \item Consultation with relevant specialists
\end{itemize}
\section{Prevention}
Treatment focuses on positive airway pressure (PAP) therapy (continuous PAP, auto-CPAP, bilevel PAP) as The first-choice treatment, mandibular advancement devices for mild-to-moderate OSA, hypoglossal nerve stimulation therapy, positional therapy, weight loss and lifestyle modification, and upper airway surgery for patients with surgically correctable anatomy or PAP intolerance.

\begin{itemize}[noitemsep]
  \item Maintain a healthy lifestyle with balanced diet and exercise
  \item Avoid known risk factors
  \item Attend regular medical check-ups for early detection
  \item Follow recommended screening guidelines
\end{itemize}
\section{Treatment Options}
Medications to manage symptoms and address underlying mechanisms Lifestyle modifications and self-care measures Physical therapy and rehabilitation Surgical intervention when indicated Regular monitoring and follow-up care Patient education and support services

\begin{itemize}[noitemsep]
  \item Medications to manage symptoms and address underlying mechanisms
  \item Lifestyle modifications and self-care measures
  \item Physical therapy and rehabilitation
  \item Surgical intervention when indicated
  \item Regular monitoring and follow-up care
\end{itemize}
\section{Living with It}
\begin{itemize}[noitemsep]
  \item Follow your treatment plan as prescribed
  \item Maintain regular follow-up with your healthcare provider
  \item Make necessary lifestyle adjustments
  \item Seek support from family, friends, or support groups
  \item Stay informed about your condition
\end{itemize}
\section{Outlook}
\begin{itemize}[noitemsep]
  \item The outlook is generally positive with treatment
  \item Untreated OSA is linked to high blood pressure, atrial fibrillation, heart failure, coronary artery disease, stroke, and type 2 diabetes.
\end{itemize}
